% ==========================================
% BAB VI EVALUASI
% ==========================================
\cleardoublepage
\chapter{EVALUASI}
\label{chap:evaluasi}
Bab ini membahas hasil pengujian dan evaluasi terhadap antarmuka aplikasi tracker ITB Ultra-Marathon yang telah diimplementasikan pada Bab V. Pengujian dilakukan melalui tiga tahap, yaitu \textit{Functionality Testing} untuk memastikan setiap fitur berjalan sesuai kebutuhan fungsional, \textit{System Integration Testing} untuk memastikan komunikasi antara front-end dan back-end berjalan dengan baik, serta \textit{User Acceptance Testing} untuk mengukur penerimaan pengguna terhadap antarmuka yang telah dikembangkan.

\section{\textit{Functionality Testing}}
\textit{Functionality Testing} dilakukan menggunakan metode \textit{black-box testing} dengan teknik \textit{equivalence partitioning}, yaitu pengujian yang berfokus pada keluaran sistem terhadap suatu masukan tanpa memeriksa struktur kode program secara langsung. Pengujian ini bertujuan untuk memastikan bahwa setiap kebutuhan fungsional (FR) yang telah didefinisikan pada Tabel III.3 dapat dijalankan dengan benar oleh antarmuka yang telah diimplementasikan, baik pada aplikasi mobile maupun aplikasi web.

Perangkat keras yang digunakan untuk pengujian \textit{website} dapat dilihat pada Tabel \ref{tab:hwt}.
\input{tables/hw-testw.tex}

Perangkat keras yang digunakan untuk pengujian aplikasi mobile dapat dilihat pada Tabel \ref{tab:hw-test}.
\begin{table}[h]
\centering
\caption{Perangkat Keras Pengujian Aplikasi Mobile}
\label{tab:hw-test}
\begin{tabular}{|p{5cm}|p{8cm}|}
\hline
\textbf{Spesifikasi} & \textbf{Detail} \\ \hline
Perangkat            & Samsung Galaxy Note10+ \\ \hline
Model                & SM-N975F/DS \\ \hline
Sistem Operasi       & Android 10 \\ \hline
Antarmuka            & One UI 2.5 \\ \hline
Versi Kernel         & 4.14.113-19542039 \\ \hline
\end{tabular}
\end{table}
Pengujian dilakukan menggunakan peramban Google Chrome untuk aplikasi web dan perangkat Android untuk aplikasi mobile. Pada pengujian aplikasi mobile, simulasi perpindahan lokasi dilakukan menggunakan aplikasi \textit{Fake GPS Simulator Joystick} untuk memverifikasi fungsi-fungsi yang bergantung pada data lokasi. Skenario pengujian disusun berdasarkan peran pengguna, kemudian dikelompokkan berdasarkan platform pengujian, yaitu aplikasi mobile dan aplikasi web. Setiap skenario dirancang untuk memverifikasi satu kebutuhan fungsional secara spesifik.

Rincian skenario serta hasil pengujian pada aplikasi mobile, yang digunakan oleh peran Pelari, disajikan pada Tabel \ref{tab:functionality-testing-mobile}.

% ==========================================
% TABEL VI.1a HASIL FUNCTIONALITY TESTING - MOBILE
% ==========================================
\begin{longtable}{|p{1.2cm}|p{4.5cm}|p{5cm}|p{2cm}|}
\caption{Hasil \textit{Functionality Testing} - Mobile (Pelari)} \label{tab:functionality-testing-mobile} \\

\hline
\textbf{ID} & \textbf{Skenario Pengujian} & \textbf{Hasil yang Diharapkan} & \textbf{Status} \\ \hline
\endfirsthead

\caption[]{Hasil \textit{Functionality Testing} - Mobile (Pelari) (lanjutan)} \\
\hline
\textbf{ID} & \textbf{Skenario Pengujian} & \textbf{Hasil yang Diharapkan} & \textbf{Status} \\ \hline
\endhead

\hline
\multicolumn{4}{r}{\textit{Bersambung ke halaman berikutnya}} \\
\endfoot

\hline
\endlastfoot

TC-01 & Pelari bergerak di sepanjang rute & ETA, jarak tersisa, waktu tempuh, dan \textit{pace} terhitung serta diperbarui secara dinamis & Berhasil \\ \hline

TC-02 & Pelari menekan tombol pusatkan kamera (\textit{center}) & Peta otomatis bergeser dan memusatkan tampilan ke titik lokasi pelari saat ini & Berhasil \\ \hline

TC-03 & Pelari menekan tombol \textit{emergency} (SOS) & Tampilan darurat muncul dan hitung mundur peringatan dimulai & Berhasil \\ \hline

TC-04 & Pelari menekan tombol menu utama & Laci menu terbuka dan menampilkan opsi navigasi tambahan & Berhasil \\ \hline

TC-05 & Pelari menekan tombol info tim & Panel informasi muncul menampilkan data tim dan anggota & Berhasil \\ \hline

TC-06 & Pelari mengubah tuas (\textit{switch}) \textit{visibility} & Status visibilitas berubah antara \textit{public} atau \textit{private} & Berhasil \\ \hline

TC-07 & Pelari menekan tombol \textit{logout} & Sesi pengguna berakhir dan aplikasi kembali ke layar \textit{login} & Berhasil \\ \hline

TC-08 & Pelari menekan layar (\textit{tap}) 2x saat kondisi \textit{emergency} & Kondisi \textit{emergency} dibatalkan (\textit{cancel}) dan kembali ke mode peta & Berhasil \\ \hline

TC-09 & Pelari menekan navigasi arah (sebelum/selanjutnya) & Peta berpindah dan memfokuskan (\textit{spectating}) pelari lain dalam tim yang sama & Berhasil \\ \hline

\end{longtable}

Sementara itu, rincian skenario serta hasil pengujian pada aplikasi web, yang mencakup peran Penonton Public, Penonton Private, Ketua Tim, Supporter, dan Panitia, disajikan pada Tabel \ref{tab:functionality-testing-web}.

% ==========================================
% TABEL VI.1b HASIL FUNCTIONALITY TESTING - WEB
% ==========================================
\begin{longtable}{|p{1.2cm}|p{4.5cm}|p{5cm}|p{2cm}|}
\caption{Hasil \textit{Functionality Testing} - Web} \label{tab:functionality-testing-web} \\

\hline
\textbf{ID} & \textbf{Skenario Pengujian} & \textbf{Hasil yang Diharapkan} & \textbf{Status} \\ \hline
\endfirsthead

\caption[]{Hasil \textit{Functionality Testing} - Web (lanjutan)} \\
\hline
\textbf{ID} & \textbf{Skenario Pengujian} & \textbf{Hasil yang Diharapkan} & \textbf{Status} \\ \hline
\endhead

\hline
\multicolumn{4}{r}{\textit{Bersambung ke halaman berikutnya}} \\
\endfoot

\hline
\endlastfoot

\multicolumn{4}{|c|}{\textbf{Penonton Public}} \\ \hline

TC-10 & Pengguna membuka halaman web publik & Peta menampilkan titik pelari-pelari yang membuka visibilitas publik & Berhasil \\ \hline

TC-11 & Pengguna mengeklik salah satu titik pelari & Muncul detail informasi pelari beserta nama timnya & Berhasil \\ \hline

TC-12 & Pengguna mencari tim dari daftar direktori tim & Peta otomatis bergeser mengarahkan tampilan ke titik tim tersebut & Berhasil \\ \hline

TC-13 & Pengguna mencari \textit{water station} dari direktori & Peta otomatis bergeser mengarahkan tampilan ke titik lokasi \textit{water station} & Berhasil \\ \hline


\multicolumn{4}{|c|}{\textbf{Ketua Tim}} \\ \hline

TC-14 & Sistem mendeteksi kondisi \textit{emergency} atau tersesat & Antarmuka Ketua menerima notifikasi masuk & Berhasil \\ \hline

TC-15 & Ketua mengeklik titik pelari pada peta & Peta menampilkan posisi pelari dan rincian detailnya & Berhasil \\ \hline

TC-16 & Ketua mengakses opsi bagikan & Tautan untuk \textit{supporter} dibuat dan siap disalin & Berhasil \\ \hline


\multicolumn{4}{|c|}{\textbf{Supporter}} \\ \hline

TC-17 & Sistem mendeteksi kondisi \textit{emergency} atau tersesat & Antarmuka Supporter menerima notifikasi peringatan & Berhasil \\ \hline

TC-18 & Supporter mengeklik titik pelari yang dipantau & Peta menampilkan posisi pelari dan rincian detailnya & Berhasil \\ \hline


\multicolumn{4}{|c|}{\textbf{Panitia}} \\ \hline

TC-19 & Panitia mengeklik titik pelari pada peta & Peta menampilkan posisi detail pelari dari seluruh tim & Berhasil \\ \hline

TC-20 & Panitia membuka daftar informasi tim & Menampilkan rincian lengkap dari tim dan anggotanya & Berhasil \\ \hline

TC-21 & Panitia membuka modul daftar pelari & Menampilkan daftar status dari masing-masing pelari & Berhasil \\ \hline

TC-22 & Panitia membuka dasbor statistik & Tampil rekapitulasi angka statistik perlombaan & Berhasil \\ \hline

TC-23 & Sistem mendeteksi adanya pelari yang SOS/tersesat & Peringatan darurat masuk ke halaman panitia dengan penanda merah & Berhasil \\ \hline

\multicolumn{4}{|c|}{\textbf{Penonton Private}} \\ \hline
TC-24 & Penonton \textit{private} membuka halaman web & Halaman menampilkan metrik aktivitas (\textit{activity metric}) dari 1 pelari yang dipantau & Berhasil \\ \hline

\end{longtable}

Berdasarkan hasil pengujian pada Tabel \ref{tab:functionality-testing-mobile} dan Tabel \ref{tab:functionality-testing-web}, seluruh 24 skenario pengujian fungsionalitas yang meliputi 9 skenario pada aplikasi mobile dan 15 skenario pada aplikasi web untuk berbagai peran pengguna berhasil dijalankan sesuai dengan hasil yang diharapkan, sehingga diperoleh tingkat keberhasilan pengujian fungsionalitas sebesar 100\%. Hal ini menunjukkan bahwa antarmuka yang dikembangkan, baik pada platform mobile maupun web, telah memenuhi seluruh spesifikasi fungsional yang dirancang pada Bab IV.

% \section{\textit{System Integration Testing}}
% \textit{System Integration Testing} dilakukan untuk memastikan bahwa komunikasi data antara front-end (aplikasi mobile berbasis React Native/Expo dan aplikasi web berbasis React) dengan back-end melalui \textit{API Gateway} dapat berjalan dengan baik. Pengujian ini berfokus pada validasi setiap titik integrasi (\textit{endpoint}) yang telah dirancang pada rancangan \textit{behavioral} (SK-01 hingga SK-08) pada Bab IV, termasuk pengujian terhadap skenario kegagalan seperti gangguan jaringan dan terputusnya sinyal GPS.

% Pengujian dilakukan menggunakan kombinasi \textit{end-to-end testing} secara manual pada aplikasi serta pengujian \textit{endpoint} API menggunakan Postman untuk memverifikasi format permintaan (\textit{request}) dan tanggapan (\textit{response}) yang dikembalikan oleh backend. Rincian skenario serta hasil pengujian integrasi disajikan pada Tabel \ref{tab:integration-testing}.

% % ==========================================
% TABEL VI.2 HASIL SYSTEM INTEGRATION TESTING
% ==========================================
\begin{longtable}{p{1.3cm} p{1.6cm} p{3cm} p{4.3cm} p{3cm} p{1.5cm}}
\caption{Hasil \textit{System Integration Testing}} \label{tab:integration-testing} \\

\hline
\textbf{ID} & \textbf{Skenario} & \textbf{Endpoint} & \textbf{Skenario Pengujian} & \textbf{Hasil yang Diharapkan} & \textbf{Status} \\ \hline
\endfirsthead

\multicolumn{6}{c}{\tablename\ \thetable{} -- lanjutan dari halaman sebelumnya} \\ \hline
\textbf{ID} & \textbf{Skenario} & \textbf{Endpoint} & \textbf{Skenario Pengujian} & \textbf{Hasil yang Diharapkan} & \textbf{Status} \\ \hline
\endhead

\hline \multicolumn{6}{r}{Bersambung ke halaman berikutnya} \\
\endfoot

\hline
\endlastfoot

IT-01 & SK-01 & \texttt{POST /location} & Front-end mobile mengirimkan koordinat, \textit{timestamp}, dan ID peserta secara berkala ke backend & Backend menyimpan data lokasi dan mengembalikan respons \texttt{200 OK} & Berhasil \\
IT-02 & SK-01 & \texttt{POST /location} & Koneksi jaringan diputus secara sengaja saat pengiriman lokasi berlangsung & Data lokasi disimpan sementara pada \textit{cache} lokal dan dikirim ulang pada interval berikutnya & Berhasil \\
IT-03 & SK-02 & \texttt{GET /eta-metrics} & Front-end meminta data ETA, ETD, dan metrik aktivitas menggunakan lokasi terkini & Backend mengembalikan nilai ETA, ETD, dan metrik aktivitas yang sesuai dengan data rute & Berhasil \\
IT-04 & SK-02 & \texttt{GET /eta-metrics} & Sinyal GPS perangkat peserta diputus sementara & Backend mengembalikan respons \texttt{204 No Content} dan antarmuka menampilkan status menunggu sinyal GPS & Berhasil \\
IT-05 & SK-03 & \texttt{POST /emergency} & Peserta menyelesaikan \textit{countdown} SOS tanpa pembatalan & Backend menyimpan status darurat peserta dan meneruskannya ke aktor terkait & Berhasil \\
IT-06 & SK-04 & \texttt{PATCH /visibility} & Peserta mengubah status visibilitas lokasi & Backend memperbarui status visibilitas dan mengembalikan respons \texttt{200 OK} beserta status terbaru & Berhasil \\
IT-07 & SK-05 & \texttt{GET /tracking} & Aktor dengan akses publik atau kode tim yang valid membuka peta pelacakan & Backend mengembalikan posisi pelari dan memperbarui data secara berkala & Berhasil \\
IT-08 & SK-05 & \texttt{GET /tracking} & Aktor mencoba mengakses data peserta dengan visibilitas privat tanpa kode tim & Backend mengembalikan respons \texttt{403 Forbidden} & Berhasil \\
IT-09 & SK-06 & \texttt{POST /invite-link} & Panitia atau Ketua Tim membuat tautan undangan & Backend mengembalikan respons \texttt{201 Created} beserta tautan undangan yang valid & Berhasil \\
IT-10 & SK-07 & \texttt{GET /emergency/\{id\}} & Panitia/Supporter menerima notifikasi darurat dan meminta detail peserta & Backend mengembalikan detail lokasi dan status peserta yang bersangkutan & Berhasil \\
IT-11 & SK-08 & \texttt{GET /runner/\{id\}/location} & Supporter menekan tombol navigasi menuju lokasi pelari & Backend mengembalikan koordinat terkini dan aplikasi membentuk URL Google Maps yang valid & Berhasil \\
IT-12 & SK-08 & \texttt{GET /runner/\{id\}/location} & Sinyal GPS pelari yang dituju terputus & Backend mengembalikan respons \texttt{204 No Content} dan aplikasi menampilkan pesan lokasi tidak tersedia & Berhasil \\
IT-13 & -- & Front-end Mobile $\leftrightarrow$ \textit{API Gateway} & Pengujian autentikasi berbasis \textit{invite link} dari aplikasi mobile menuju backend & Token akses dan \textit{refresh token} diterbitkan dan diterima dengan benar oleh aplikasi mobile & Berhasil \\
IT-14 & -- & Front-end Website $\leftrightarrow$ \textit{API Gateway} & Panitia mengelola data rute dan checkpoint melalui antarmuka web & Perubahan data tersimpan pada PostgreSQL dan tersinkronisasi pada dashboard pemantauan & Berhasil \\

\end{longtable}


% Berdasarkan hasil pengujian pada Tabel \ref{tab:integration-testing}, seluruh skenario integrasi antara front-end dan back-end berhasil dijalankan sesuai dengan hasil yang diharapkan. Mekanisme penanganan kegagalan, seperti penyimpanan data lokasi sementara pada \textit{cache} lokal saat jaringan terputus (IT-02) serta pembatasan akses data peserta dengan visibilitas privat (IT-08), juga terbukti berfungsi dengan baik. Dengan demikian, dapat disimpulkan bahwa integrasi antara front-end dan back-end pada aplikasi tracker ITB Ultra-Marathon telah berjalan secara efektif dan sesuai dengan rancangan komunikasi data yang telah ditetapkan.

\section{\textit{User Acceptance Testing}}
\textit{User Acceptance Testing} (UAT) dilakukan untuk mengukur tingkat penerimaan pengguna terhadap antarmuka aplikasi \textit{tracker} ITB Ultra-Marathon yang telah dikembangkan, khususnya pada antarmuka \textit{Public}. Pengujian ini dilakukan dengan menyebarkan tautan aplikasi untuk disimulasikan secara mandiri, yang diikuti oleh 35 responden menggunakan berbagai jenis perangkat keras (komputer maupun \textit{smartphone}). Rincian demografi responden berdasarkan perangkat yang digunakan disajikan pada Tabel \ref{tab:uat-demografi}.

% ==========================================
% TABEL VI.3 DEMOGRAFI RESPONDEN UAT
% ==========================================
\begin{table}[H]
\centering
\caption{Demografi Responden \textit{User Acceptance Testing}}
\label{tab:uat-demografi}
\begin{tabular}{lll}
\hline
\textbf{Kategori} & \textbf{Kelompok} & \textbf{Jumlah Responden} \\ \hline
Peran Pengguna & Peserta (Pelari)                  & 8  \\
               & Ketua Tim                         & 3  \\
               & Supporter                         & 5  \\
               & Panitia                           & 2  \\
               & Penonton/\textit{Spectator}       & 2  \\ \hline
Pengalaman ITB Ultra-Marathon & Pernah mengikuti/berpartisipasi & 14 \\
                               & Belum pernah berpartisipasi     & 6  \\ \hline
\textbf{Total Responden} & & \textbf{20} \\ \hline
\end{tabular}
\end{table}


Setelah melakukan simulasi penggunaan aplikasi, responden diminta untuk mengisi kuesioner yang disusun menggunakan skala Likert 1 (sangat tidak setuju) hingga 5 (sangat setuju). Terdapat 26 pernyataan pada kuesioner yang disusun untuk mengukur berbagai aspek utama, yaitu kemudahan penggunaan, kejelasan antarmuka, persepsi kinerja, kualitas informasi, serta kepuasan pengguna secara keseluruhan. Hasil rekapitulasi nilai rata-rata dari kuesioner UAT tersebut disajikan pada Tabel \ref{tab:uat-hasil}.

% ==========================================
% TABEL VI.4 HASIL KUESIONER USER ACCEPTANCE TESTING
% ==========================================
\begin{longtable}{|p{1.2cm}|p{8.5cm}|p{2cm}|}
\caption{Hasil Kuesioner \textit{User Acceptance Testing}} \label{tab:uat-hasil} \\

\hline
\textbf{Kode} & \textbf{Pernyataan} & \textbf{Rata-rata} \\ \hline
\endfirsthead

\caption[]{Hasil Kuesioner \textit{User Acceptance Testing} (lanjutan)} \\
\hline
\textbf{Kode} & \textbf{Pernyataan} & \textbf{Rata-rata} \\ \hline
\endhead

\hline
\multicolumn{3}{r}{\textit{Bersambung ke halaman berikutnya}} \\
\endfoot

\hline
\endlastfoot

\multicolumn{3}{|c|}{\textbf{Kemudahan Penggunaan}} \\ \hline
P01 & Saya dapat memahami tujuan halaman dalam beberapa detik pertama. & 4,23 \\ \hline
P02 & Informasi ringkasan (Total Tim, Active, Finished) mudah ditemukan. & 4,74 \\ \hline
P10 & Navigasi peta (geser dan zoom) terasa nyaman digunakan. & 4,63 \\ \hline
P12 & Saya dapat menemukan cara membuka Detail Pelari dengan mudah. & 4,34 \\ \hline
P16 & Saya dapat menutup panel dan kembali menggunakan peta dengan mudah. & 4,69 \\ \hline
P17 & Menu Direktori mudah ditemukan. & 4,60 \\ \hline
P18 & Saya dapat menemukan tim yang dipilih dengan cepat. & 4,66 \\ \hline
P20 & Saya dapat menemukan lokasi yang dipilih dengan cepat. & 4,60 \\ \hline
P22 & Halaman \textit{Public Tracking} mudah digunakan. & 4,54 \\ \hline

\multicolumn{3}{|c|}{\textbf{Kejelasan Antarmuka}} \\ \hline
P04 & Tata letak halaman membantu saya menemukan informasi penting dengan cepat. & 4,57 \\ \hline
P05 & Ukuran tulisan dan elemen nyaman dilihat pada perangkat saya. & 4,69 \\ \hline
P06 & Tampilan warna dan mode gelap mendukung keterbacaan halaman. & 4,51 \\ \hline
P08 & Garis rute perlombaan terlihat dengan jelas. & 4,57 \\ \hline
P09 & Saya dapat membedakan elemen yang tampil pada peta. & 4,46 \\ \hline
P14 & Tata letak informasi pada Detail Pelari tersusun dengan baik. & 4,54 \\ \hline
P15 & Ukuran panel Detail Pelari tidak mengganggu penggunaan peta. & 4,43 \\ \hline

\multicolumn{3}{|c|}{\textbf{Persepsi Kinerja}} \\ \hline
P03 & Perpindahan (slide) pada carousel sponsor berjalan dengan baik. & 4,37 \\ \hline
P07 & Peta berhasil dimuat dengan baik. & 4,57 \\ \hline
P11 & Performa peta tetap lancar saat digunakan. & 4,57 \\ \hline
P19 & Peta berhasil berpindah ke titik tim yang sesuai. & 4,66 \\ \hline
P21 & Peta berhasil berpindah ke titik lokasi yang sesuai. & 4,54 \\ \hline

\multicolumn{3}{|c|}{\textbf{Kualitas Informasi}} \\ \hline
P13 & Informasi pada Detail Pelari (progress, leg, distance) ditampilkan dengan lengkap. & 4,60 \\ \hline
P23 & Informasi yang tersedia sudah cukup untuk melakukan pemantauan. & 4,69 \\ \hline

\multicolumn{3}{|c|}{\textbf{Kepuasan Pengguna}} \\ \hline
P24 & Saya merasa nyaman menggunakan halaman ini. & 4,51 \\ \hline
P25 & Saya bersedia menggunakan halaman ini kembali apabila tersedia pada acara sebenarnya. & 4,69 \\ \hline
P26 & Secara keseluruhan saya puas terhadap pengalaman penggunaan halaman \textit{Public Tracking}. & 4,49 \\ \hline

\multicolumn{2}{|r|}{\textbf{Rata-rata Keseluruhan}} & \textbf{4,56} \\ \hline

\end{longtable}


Berdasarkan hasil pada Tabel \ref{tab:uat-hasil}, diperoleh rata-rata skor keseluruhan sebesar 4,56 dari skala maksimum 5,0. Skor tertinggi diperoleh pada aspek kemudahan menemukan informasi ringkasan (P02) dengan nilai rata-rata 4,74, serta aspek kesediaan menggunakan aplikasi ini kembali di masa mendatang (P25) dengan nilai 4,69. Hal ini menunjukkan tingkat kepuasan yang sangat tinggi dan penerimaan yang positif dari pengguna. Sementara itu, skor terendah berada pada pernyataan mengenai kemudahan memahami tujuan halaman pada pandangan pertama (P01) dengan nilai 4,23. Meskipun merupakan skor terendah, nilai tersebut masih berada pada rentang kategori “setuju” (di atas 4,00) yang mengindikasikan bahwa halaman pada akhirnya tetap dapat dipahami dengan baik oleh pengguna.

Secara keseluruhan, hasil UAT menunjukkan bahwa antarmuka \textit{Public} ITB Ultra-Marathon yang dikembangkan telah memenuhi kriteria keberhasilan dari sisi kemudahan penggunaan, kejelasan penyajian informasi visual peta, persepsi kelancaran kinerja, serta tingkat kepuasan pengguna akhir yang tinggi.